\section{相关工作}

本文的技术核心来自三条相对独立但可在同一形式体系中对齐的传统：空间填充曲线与多尺度局部性、图同构与 Weisfeiler--Leman（WL）细化、以及约束满足（CSP）与受约束搜索复杂度。

\subsection{空间填充曲线与局部性组织}

Hilbert 曲线是经典的空间填充曲线之一，其递归结构提供了天然的多尺度块分解，同时具有良好的局部性保持性质：沿曲线相邻的索引往往对应几何空间中的近邻。系统性介绍可见 \cite{sagan1994space}。本文将 Hilbert 曲线族视为观察者的“尺子”：它并不改变被研究对象的集合本身，而是改变观察者组织局部性、选择扫描顺序与多尺度块结构的方式。

\subsection{WL 细化与“塌缩速度”}

Weisfeiler--Leman 细化（1-WL 即颜色细化）给出一种可审计的迭代过程：每轮把节点颜色替换为其旧颜色与邻居颜色多重集的组合，并进行确定性压缩 \cite{weisfeiler1968reduction}。本文把“解析到唯一微观区分所需的迭代步数”作为一类观察者时间指标，从而把“塌缩速度”翻译为可复现实验量。

\subsection{约束满足与观察者时间}

在观察者只见投影信息的情况下，“回到过去/向上恢复”往往变成一个受约束搜索问题。该领域的标准语言是约束满足与推理/回溯框架 \cite{rossi2006csp}。本文的观点是：时空结构可被定义为搜索过程生成的二级图对象，且其复杂度（访问节点数、回溯次数、传播停滞等）构成“观察者时间”的纯数学刻画。

\begin{remark}
本文进一步把时间轴的本体成分落实为 uplift/历史胶带的编码长度（负信息时间 \(\tautime\)），并在最小假设下给出 \(\kappaobs\sim\tautime\) 的同阶等价关系；因此，上述“复杂度指标”也可被等价地理解为对负信息时间的可审计 proxy。
\end{remark}

\subsection{与元胞自动机与信息组织的关系}

在复杂系统与元胞自动机传统中，局部规则诱导的宏观结构、以及“可见层/隐变量层”的信息组织是重要主题（参见综述性讨论 \cite{wolfram2002newkind}）。本文不以物理解释作为推理前提，而是强调：当本体被抽象为约束系统时，观察者的可见化与推断策略本身即可产生复杂的时空图与时间尺度差异。

